\subsubsection{Complete Example for \texorpdfstring{$w = 101$}{w = 101} (Exam Question)}\label{sec:complete-example-for-w101-exam-question}

Having understood the construction of the oracle (\autopageref{sec:quantum-oracle-construction}) and the role of phase kickback (\autopageref{sec:phase-kickback-intuition}), we can now \hl{execute the entire Bernstein-Vazirani algorithm} for the hidden bitstring:
\begin{equation*}
    w = 101
\end{equation*}
Rather than focusing on the mathematical derivation, we simply follow the quantum state after each layer of the circuit.

\highspace
\begin{deepeningbox}[: Quantum Computing Exam Question]
    \textcolor{Red2}{\faIcon{exclamation-triangle}} Note that this is one of the \hl{exam questions for the course}, so it is important to understand and be able to reproduce the steps. The professor will generally ask us to draw the circuit for the Bernstein-Vazirani algorithm and then execute it for a specific hidden bit string w. Another common option is to provide the oracle circuit directly and ask us to draw the complete Bernstein-Vazirani algorithm circuit and execute it for the specific hidden bitstring $w$ obtained from the oracle circuit.

    The following is an example of such a question, taken from the June 12, 2026 exam. We will not provide the solution here, but we will go through a complete, worked example for w = 101 in the following sections. Understanding the w = 101 example will enable us to answer exam questions involving any hidden bitstring w, as with other exam questions concerning the Bernstein-Vazirani algorithm.
    \begin{examplebox}[: Exercise taken from the 2026 exam, June 12th]
        Bernstein-Vazirani Algorithm (8 points) exercise.
        \begin{itemize}
            \item Briefly describe the problem solved by the Bernstein-Vazirani algorithm and draw the circuit of the algorithm in its general form.

            \item Consider then the oracle $U_f$ implemented by the quantum circuit:

            \begin{center}
                \begin{quantikz}
                    \lstick{Data: $q_1$}        & \targ{}   & \ctrl{2} & \targ{}   & \ctrl{2} & \\
                    \lstick{$q_2$}              & \ctrl{-1} &          & \ctrl{-1} &          & \\
                    \lstick{Computation: $q_3$} &           & \targ{}  &           & \targ{}  &
                \end{quantikz}
            \end{center}

            \item Draw the full circuit of the Bernstein-Vazirani algorithm applied to the oracle $U_f$.

            \item Apply the Bernstein-Vazirani algorithm to the oracle $U_f$ and determine the hidden bitstring encoded in the oracle.
        \end{itemize}
    \end{examplebox}
\end{deepeningbox}

\highspace
\begin{flushleft}
    \textcolor{Green3}{\faIcon{play-circle} \textbf{Step 1: Initial State}}
\end{flushleft}
The computation starts with three data qubits initialized to \ket{0} and one ancilla qubit initialized to \ket{1}:
\begin{equation*}
    \ket{\psi_0} = \ket{000}\ket{1}
\end{equation*}
The data register (i.e., the first three qubits) contains no information about the hidden bitstring $w$.
\begin{center}
    \begin{quantikz}
            \lstick{$\ket{0}$}\slice{$\ket{\psi_0}$}  & \\
            \lstick{$\ket{0}$}                        & \\
            \lstick{$\ket{0}$}                        & \\
            \lstick{$\ket{1}$}                        &
        \end{quantikz}
\end{center}

\highspace
\begin{flushleft}
    \textcolor{Green3}{\faIcon{project-diagram} \textbf{Step 2: First Hadamard Layer}}
\end{flushleft}
Hadamard gates are applied to every qubit. The data register becomes:
\begin{equation*}
    \ket{0} \to \ket{+} = \frac{1}{\sqrt{2}}(\ket{0} + \ket{1})
\end{equation*}
While the ancilla qubit becomes:
\begin{equation*}
    \ket{1} \to \ket{-} = \frac{1}{\sqrt{2}}(\ket{0} - \ket{1})
\end{equation*}
The quantum state is therefore:
\begin{equation*}
    \ket{\psi_1} = \ket{+,+,+}\ket{-}
\end{equation*}
At this point, \textbf{every data qubit is in an equal superposition}, and the \textbf{ancilla is prepared for phase kickback}.
\begin{center}
    \begin{quantikz}
            \lstick{$\ket{0}$}\slice{$\ket{\psi_0}$}  & \gate{H}\slice{$\ket{\psi_1}$}  & \\
            \lstick{$\ket{0}$}                        & \gate{H}                        & \\
            \lstick{$\ket{0}$}                        & \gate{H}                        & \\
            \lstick{$\ket{1}$}                        & \gate{H}                        &
        \end{quantikz}
\end{center}

\newpage

\begin{flushleft}
    \textcolor{Green3}{\faIcon{magic} \textbf{Step 3: Oracle Application}}
\end{flushleft}
The hidden string is:
\begin{equation*}
    w = 101
\end{equation*}
So the oracle contains two CNOT gates:
\begin{equation*}
    f_w (x) = \left(1 \cdot x_1\right) \oplus \left(0 \cdot x_2\right) \oplus \left(1 \cdot x_3\right) = x_1 \oplus x_3
\end{equation*}
One controlled by the first qubit $x_1$ and one controlled by the third qubit $x_3$. Because the ancilla is in the state \ket{-}, these CNOTs do not flip the ancilla. INstead, they produce \textbf{phase kickback}.
\begin{center}
    \begin{quantikz}
        \lstick{$\ket{0}$}\slice{$\ket{\psi_0}$}  & \gate{H}\slice{$\ket{\psi_1}$}  &  & \ctrl{3}\gategroup[4,steps=3,style={inner sep=3pt}]{$U_f$}  &           &           &  \slice{$\ket{\psi_2}$}  &  \\
        \lstick{$\ket{0}$}                        & \gate{H}                        &  &                                                             & \gate{I}  &           &                          &  \\
        \lstick{$\ket{0}$}                        & \gate{H}                        &  &                                                             &           & \ctrl{1}  &                          &  \\
        \lstick{$\ket{1}$}                        & \gate{H}                        &  & \targ{}                                                     &           & \targ{}   &                          & 
    \end{quantikz}
\end{center}

\noindent
Without providing the full mathematical derivation as we did on \autopageref{sec:phase-kickback-intuition}, we can take some shortcuts. A data qubit in the state \ket{+} becomes \ket{-} if it is a control for a CNOT gate and remains \ket{+} if it is not. Thus, the first and third qubits become \ket{-}, while the second qubit remains \ket{+}. Due to phase kickback, the ancilla remains in the state \ket{-}.
\begin{itemize}
    \item The first qubit changes from \ket{+} to \ket{-} because it is a control for a CNOT gate.
    \item The second qubit remains in the state \ket{+} because it is not a control for any CNOT gate.
    \item The third qubit changes from \ket{+} to \ket{-} because it is a control for a CNOT gate.
    \item The ancilla remains in the state \ket{-} due to phase kickback
\end{itemize}
The resulting quantum state is therefore:
\begin{equation*}
    \ket{\psi_2} = \ket{-,+,-}\ket{-}
\end{equation*}
Notice that the pattern of plus and minus states now exactly matches the hidden string:
\begin{equation*}
    w = 101 \quad \iff \quad \left(\ket{-}, \ket{+}, \ket{-}\right)
\end{equation*}
The oracle has encoded the hidden bits into the phases of the data qubits.

\newpage

\begin{flushleft}
    \textcolor{Green3}{\faIcon{compress-arrows-alt} \textbf{Step 4: Final Hadamard Layer}}
\end{flushleft}
The final Hadamard gates convert the phase states back into computational basis states. Using:
\begin{equation*}
    H\ket{+} = \ket{0}, \quad H\ket{-} = \ket{1}
\end{equation*}
We obtain:
\begin{equation*}
    \ket{-} \to \ket{1}, \quad \ket{+} \to \ket{0}, \quad \ket{-} \to \ket{1}
\end{equation*}
Therefore:
\begin{equation*}
    \ket{\psi_3} = \ket{101}\ket{-}
\end{equation*}
The first three qubits now contain exactly the hidden bitstring $w = 101$, while the ancilla remains in the state \ket{-}.
\begin{center}
    \begin{quantikz}
        \lstick{$\ket{0}$}\slice{$\ket{\psi_0}$}  & \gate{H}\slice{$\ket{\psi_1}$}  &  & \ctrl{3}\gategroup[4,steps=3,style={inner sep=3pt}]{$U_f$}  &           &           &  \slice{$\ket{\psi_2}$}  & \gate{H}\slice{$\ket{\psi_3}$}  & \\
        \lstick{$\ket{0}$}                        & \gate{H}                        &  &                                                             & \gate{I}  &           &                          & \gate{H}                        & \\
        \lstick{$\ket{0}$}                        & \gate{H}                        &  &                                                             &           & \ctrl{1}  &                          & \gate{H}                        & \\
        \lstick{$\ket{1}$}                        & \gate{H}                        &  & \targ{}                                                     &           & \targ{}   &                          &                                 &
    \end{quantikz}
\end{center}

\highspace
\begin{flushleft}
    \textcolor{Green3}{\faIcon{search} \textbf{Step 5: Measurement}}
\end{flushleft}
Finally, the data register is measured. Since its state is already \ket{101}, the measurement outcome is \textbf{deterministically} (Bernstein-Vazirani is a deterministic algorithm):
\begin{equation*}
    \text{Measurement outcome: } 101 \impliedby P(101) = 1
\end{equation*}
The ancilla is not measured because it contains no useful information about the solution.
\begin{center}
    \begin{quantikz}
        \lstick{$\ket{0}$}\slice{$\ket{\psi_0}$}  & \gate{H}\slice{$\ket{\psi_1}$}  &  & \ctrl{3}\gategroup[4,steps=3,style={inner sep=3pt}]{$U_f$}  &           &           &  \slice{$\ket{\psi_2}$}  & \gate{H}\slice{$\ket{\psi_3}$}  & \meter{} \\
        \lstick{$\ket{0}$}                        & \gate{H}                        &  &                                                             & \gate{I}  &           &                          & \gate{H}                        & \meter{} \\
        \lstick{$\ket{0}$}                        & \gate{H}                        &  &                                                             &           & \ctrl{1}  &                          & \gate{H}                        & \meter{} \\
        \lstick{$\ket{1}$}                        & \gate{H}                        &  & \targ{}                                                     &           & \targ{}   &                          &                                 &
    \end{quantikz}
\end{center}

\newpage

\begin{flushleft}
    \textcolor{Red2}{\faIcon{exclamation-triangle} \textbf{What happens when the oracle contains gates that affect data qubits, but not ancillas?}}
\end{flushleft}
Suppose the oracle provided by a \emph{mysterious company} contains a quantum circuit that differs from the classic one we have seen. For instance, it \hl{may contain a CNOT gate controlled by a data qubit that targets another data qubit.} How would we show each step of the Bernstein-Vazirani algorithm in this case? What should the state of the data qubits be after applying the oracle? These questions arise from a similar exercise that was on the \textbf{June 12, 2026 exam}.

\highspace
In the previous sections, we always built the oracle starting from the hidden string:
\begin{equation*}
    f_w(x) = w \cdot x \mod 2
\end{equation*}
Which naturally leads to \textbf{CNOTs from the data qubits to the ancilla}, because the ancilla stores:
\begin{equation*}
    y \to y \oplus f_w(x)
\end{equation*}
That is only \textbf{one particular implementation of the oracle}. The oracle can be implemented in many different ways, as long as it satisfies the condition:
\begin{equation*}
    U_f \ket{x}\ket{y} = \ket{x}\ket{y \oplus f_w(x)}
\end{equation*}

\highspace
Suppose the oracle circuit is:
\begin{center}
    \begin{quantikz}
        \lstick{$q_1$}             & \targ{}   & \ctrl{2} & \targ{}   & \ctrl{2} & \\
        \lstick{$q_2$}             & \ctrl{-1} &          & \ctrl{-1} &          & \\
        \lstick{$\text{anc}=q_3$}  &           & \targ{}  &           & \targ{}  &
    \end{quantikz}
\end{center}
\noindent
Where:
\begin{itemize}
    \item $q_1$ and $q_2$ are \textbf{data qubits} (data register);
    \item $q_3$ is the \textbf{computation (ancilla) qubit}.
\end{itemize}
Notice that there are \textbf{four gates}:
\begin{enumerate}
    \item $\mathrm{CNOT } q_2 \to q_1$: A CNOT gate controlled by $q_2$ and targeting $q_1$ (data qubit);
    \item $\mathrm{CNOT } q_1 \to q_3$: A CNOT gate controlled by $q_1$ and targeting the ancilla $q_3$;
    \item $\mathrm{CNOT } q_2 \to q_1$: A CNOT gate controlled by $q_2$ and targeting $q_1$ (data qubit);
    \item $\mathrm{CNOT } q_1 \to q_3$: A CNOT gate controlled by $q_1$ and targeting the ancilla $q_3$.
\end{enumerate}
The important thing is that the two CNOTs on the first line modify the \textbf{data register itself} before it is used to control the CNOTs that target the ancilla.

\highspace
\textcolor{Green3}{\faIcon{question-circle} \textbf{Why modify the data qubits?}} Because the function may depend on a \textbf{transformed version} of the input. Instead of computing simply:
\begin{equation*}
    f(x) = x_1
\end{equation*}
The circuit may compute:
\begin{equation*}
    f\left(x_1 \oplus x_2, x_2\right)
\end{equation*}
Or some equivalent Boolean expression. The first pair of CNOTs changes the variables that are later used to control the ancilla. So the oracle is implementing a function that depends on a \textbf{transformed version of the input}, rather than the input itself:
\begin{equation*}
    U_f \ket{x, y} = \ket{g(x)} \ket{y \oplus f(g(x))}
\end{equation*}
Where $g(x)$ is the reversible transformation performed on the data qubits.

\highspace
\textcolor{Green3}{\faIcon{question-circle} \textbf{How do we recover the hidden string?}} After observing the oracle circuit, the clever part begins. Unlike the previous example with naïve oracle circuits, where the hidden bitstring was decoded by simply counting the CNOTs connected to the ancilla, we must now \textbf{determine the Boolean function implemented by the circuit}.
\begin{enumerate}
    \item \textbf{Initial data bits}. The data qubits are:
    \begin{equation*}
        q_1 = x_1, \quad q_2 = x_2, \quad q_3 = y
    \end{equation*}
    \begin{center}
    \begin{quantikz}
            \lstick{$q_1$}             & \\
            \lstick{$q_2$}             & \\
            \lstick{$\text{anc}=q_3$}  &
        \end{quantikz}
    \end{center}

    \item \textbf{First CNOT}. The first CNOT gate is controlled by $q_2$ and targets $q_1$. So we write the XOR operation:
    \begin{equation*}
        \text{CNOT } q_2 \to q_1 \implies q_1 = x_1 \oplus x_2
    \end{equation*}
    \textcolor{Green3}{\faIcon{question-circle} \textbf{Why XOR?}} Because the CNOT gate flips the target qubit if the control qubit is 1. To replicate this behavior, we use the logical XOR operation, which flips the target bit if the control bit is 1:
    \begin{center}
        \begin{tabular}{@{} c c c @{}}
            \toprule
            \textbf{Control} & \textbf{Target} & \textbf{XOR Result} \\
            \midrule
            0 & 0 & 0 \\
            0 & 1 & 1 \\
            1 & 0 & 1 \\
            1 & 1 & 0 \\
            \bottomrule
        \end{tabular}
    \end{center}
    \begin{center}
        \begin{quantikz}
            \lstick{$q_1$}             & \targ{}   & \\
            \lstick{$q_2$}             & \ctrl{-1} & \\
            \lstick{$\text{anc}=q_3$}  &           &
        \end{quantikz}
    \end{center}

    \item \textbf{First ancilla CNOT}. The next CNOT gate is controlled by $q_1$ and targets the ancilla $q_3$. So the ancilla receives $q_1$:
    \begin{equation*}
        \text{CNOT } q_1 \to q_3 \implies q_3 = y \oplus q_1 = y \oplus \left(x_1 \oplus x_2\right)
    \end{equation*}
    \begin{center}
        \begin{quantikz}
            \lstick{$q_1$}             & \targ{}   & \ctrl{2} & \\
            \lstick{$q_2$}             & \ctrl{-1} &          & \\
            \lstick{$\text{anc}=q_3$}  &           & \targ{}  &
        \end{quantikz}
    \end{center}

    \item \textbf{Second CNOT}. The next CNOT gate is controlled by $q_2$ and targets $q_1$, so it affects only the data qubits. The new value of $q_1$ is:
    \begin{align*}
        \text{CNOT } q_2 \to q_1 \implies q_1 &= \left(x_1 \oplus x_2\right) \oplus x_2 \\
        &= x_1 \oplus x_2 \oplus x_2 \\
        &\downarrow \text{XOR property: } a \oplus a = 0, \quad \forall a \in \{0, 1\} \\
        &= x_1 \oplus 0 \\
        &\downarrow \text{XOR property: } a \oplus 0 = a, \quad \forall a \in \{0, 1\} \\
        &= x_1
    \end{align*}
    So the second CNOT \textbf{undoes the effect of the first CNOT}, and $q_1$ returns to its original value $x_1$. This is a common reversible-computing trick: we can use a CNOT to modify a qubit and then use another CNOT to restore it to its original value. The \textbf{ancilla is unaffected} by this second CNOT, so it remains:
    \begin{equation*}
        q_3 = y \oplus \left(x_1 \oplus x_2\right)
    \end{equation*}
    \begin{center}
    \begin{quantikz}
            \lstick{$q_1$}             & \targ{}   & \ctrl{2} & \targ{}   & \\
            \lstick{$q_2$}             & \ctrl{-1} &          & \ctrl{-1} & \\
            \lstick{$\text{anc}=q_3$}  &           & \targ{}  &           &
        \end{quantikz}
    \end{center}

    \item \textbf{Final ancilla CNOT}. The last CNOT gate is controlled by $q_1$ and targets the ancilla $q_3$. So the ancilla receives $q_1$:
    \begin{align*}
        \text{CNOT } q_1 \to q_3 \implies q_3 &= \left(y \oplus \left(x_1 \oplus x_2\right)\right) \oplus x_1 \\
        &= y \oplus x_1 \oplus x_2 \oplus x_1 \\
        &= y \oplus x_2
    \end{align*}
    \begin{center}
    \begin{quantikz}
            \lstick{$q_1$}             & \targ{}   & \ctrl{2} & \targ{}   & \ctrl{2} & \\
            \lstick{$q_2$}             & \ctrl{-1} &          & \ctrl{-1} &          & \\
            \lstick{$\text{anc}=q_3$}  &           & \targ{}  &           & \targ{}  &
        \end{quantikz}
    \end{center}
\end{enumerate}
Therefore, the oracle implements the function:
\begin{equation*}
    f(x_1, x_2) = x_2
\end{equation*}
Which means that the hidden string is:
\begin{equation*}
    w = 01
\end{equation*}

\begin{takeawaysbox}
    If the oracle contains gates that affect also the data qubits:
    \begin{enumerate}
        \item Do \textbf{not} try to read the hidden bitstring $w$ directly from the CNOTs.
        \item Trace the Boolean value carried by every data qubit after each gate in the oracle circuit.
        \item Write the expression that reaches the ancilla.
        \item Simplify the XOR expression.
        \item Rewrite the result as:
        \begin{equation*}
            f_w(x) = w \cdot x \mod 2
        \end{equation*}
        From which we immediately read the hidden bitstring $w$.
    \end{enumerate}
\end{takeawaysbox}